\section{Route only when the host is much more expensive, caching included}\label{sec:routing}

\textbf{Claim.} Routing a session to a cheaper model saves money only when the host is much more expensive once you
count what a session actually buys (mostly cache reads and writes) and the extra requests the cheaper model makes.
On Fable, routing alone cost \CfFableSonnetArm\X\ the plain host and the full preregistered package
\SHOneRatio\X; on Opus it cost \SHFiveRatio--\CfOpusSonnetArm\X. Both directions replicated on fresh scenarios.

\paragraph{Why we measured whole sessions} Per-token list prices misled. Opus's input price is above Sonnet's, yet
an agent session is mostly cached context, and Sonnet's cache reads cost more than Opus's. The only honest test was
to run the same task twice, once on the plain host and once routed, and compare the bills.

\paragraph{Evidence} \Cref{tab:perrequest} prices one average request of the plain-host sessions (\RqMixRead\
cache-read, \RqMixWrite\ cache-write and \RqMixOut\ output tokens, measured over \RqMixAnchors\ sessions) at each
model's list prices. Sonnet costs \RqSonnetOverFable\X\ Fable per request but \RqSonnetOverOpus\X\ Opus. Sonnet
sessions also made more requests (\RpMultFable\X\ on Fable's tasks, \RpMultOpus\X\ on Opus's), so the predicted
routing ratios are \RqPredFable\ and \RqPredOpus. The measured ratios agree in sign and size (\cref{fig:replication}).
In the main campaign plain Sonnet against plain Fable, both at default effort, cost \CfFableSonnetArm\X\
(\CfFableSonnetArmCIThree), the cleanest model-only contrast; against plain Opus it cost \CfOpusSonnetArm\X\
(\CfOpusSonnetArmCIThree) \prereg\ (all pairs; the quality-filtered pair reading is in the supplement). On \SOnePanelN\ fresh scenarios the preregistered Fable configuration (routing plus medium
effort) cost \SHOneRatio\X\ plain Fable (\SHOneCI) and routing on Opus cost \SHFiveRatio\X\ (\SHFiveCI) \prereg.
Holding token counts fixed, routing on Opus would only break even if Opus's cache reads cost about
\$\BEStickyPair\ per million tokens instead of \$\PriceOpusRead\ \expl.

\begin{table}[tbp]
\centering\small
\caption{Sonnet is half Fable's price per request but Opus's equal. Cost of one average plain-host request at list
prices, its composition, and the routing ratio predicted from the per-request price and the measured request
multiplier \expl.}
\label{tab:perrequest}
\begin{tabular*}{\textwidth}{@{\extracolsep{\fill}}l S[table-format=1.4] S[table-format=2.0] S[table-format=2.0] S[table-format=2.0] c c c@{}}
\toprule
 & {\$ per request} & \multicolumn{3}{c}{share of request cost (\%)} & \multicolumn{3}{c}{routing to Sonnet} \\
\cmidrule(lr){3-5}\cmidrule(l){6-8}
Model & & {write} & {read} & {output} & per request & requests & predicted \\
\midrule
Fable 5.1 (host) & 0.1164 & 58 & 19 & 23 & 0.47 & 1.11 & 0.524 \\
Opus 5.5 (host) & 0.0554 & 49 & 32 & 20 & 0.99 & 1.38 & 1.365 \\
Sonnet 5 (routing target) & 0.0548 & 37 & 48 & 15 & -- & -- & -- \\
\bottomrule
\end{tabular*}

\end{table}

\begin{figure}[tbp]
\centering
\pgfplotslegendfromname{cumlegend}\par\smallskip
% Routing cost ratio per host across the studies (replication view).
\begin{tikzpicture}
\begin{axis}[benchmark, xmode=log, log ticks with fixed point, xtick={0.5,0.6,0.8,1,1.2,1.5,1.8}, xmin=0.42, xmax=1.95,
  scale only axis, width=8.5cm, height=5.6cm, xmajorgrids, ymin=-0.6, ymax=9.0, ytick={\CumTicks},
  yticklabels/.expanded={\CumLabels}, y tick label style={font=\footnotesize},
  xlabel={Cost ratio, routed to Sonnet / plain host (log)},  legend to name=cumlegend, legend columns=1, legend cell align=left]
\draw[black!70, line width=0.8pt] (axis cs:1,-0.6) -- (axis cs:1,8.6);
\addplot[nodes near coords, point meta=explicit symbolic, nodes near coords style={font=\scriptsize, fill=white, fill opacity=0.85, text opacity=1, inner sep=1pt, anchor=south, yshift=2.5pt, text=black!80}, only marks, mark=*, mark size=2.4pt, color=okblue, error bars/.cd, x dir=both, x explicit,
  error bar style={line width=1pt}] table[meta=v, x=r, y=y, x error minus=rm, x error plus=rp]{generated/data/cumulative.dat};
\addlegendentry{geometric-mean cost ratio, 95\,\% scenario-cluster CI}
\end{axis}
\end{tikzpicture}

\caption{The same direction in every study. Routing to Sonnet against the plain host, per host: the pilot and the
main campaign's training scenarios \expl, its test scenarios (sticky arm, all pairs) and the preregistered holdout
study (the frozen Fable configuration; Sonnet on Opus).}
\label{fig:replication}
\howtoread{Each dot is a geometric-mean cost ratio with its 95\,\% interval on a log scale; left of 1 saves money.
Every Fable row sits near one half and every Opus row above 1; the small Opus pilot is the highest.}
\end{figure}

Because two hosts and one cheap model were measured, ``much more expensive'' is a modelled threshold: no host lay
between Opus's one-to-one and Fable's two-to-one per-request price against Sonnet.

\begin{keyidea}[title={For practitioners}]
Before routing, price an average request of your own sessions on both models, including cache reads and writes, and
multiply by how many more requests the cheaper model makes. Route only when the result is clearly below one; a price
gate that does exactly this predicted both of our directions.
\end{keyidea}

Details: Supplement Sections~\ref{S-sec:campaign} and~\ref{S-sec:s1}.
